\scanpage{270}


(163)\quad Il y a une \deleted{\unsure{équivalence|détermination}}\inserted{équivalence} entre les données $(S,H,Z,\unsure{\Pi|\pi},s_0,D_1)$ d'un multigraphe à homothéties avec groupe de \inserted{$Z$-}autom $H$ \inserted{simplem.} transitif sur les $\overrightarrow A_\lambda$, un sommet marqué $s_0$ et un rayon $D_1$ incident à $Zs_0=D_0$ marqué, \inserted{avec $D_0\ne D_1$ ($\Leftrightarrow$ cas ``non dégénéré'')} d'une part,

et d'autre part : les systèmes $(H,U_0,T_H,N_H)$ d'un groupe $H$, et de ss-groupes $U_0,T_H,N_H$, satisfaisant :

a) $T_H$ normalise $U_0$, et $T_H\cap U_0=\{1\}$ (d'où produit 1/2 direct $B_H=T_H.U_0$)

b) $N_H\supset\deleted{\unsure{B_H}}T_H$, $T_H$ d'indice \deleted{\illegible{1}} 2 dans $N_H$ (qui normalise donc $T_H$) -- de sorte que les $\sigma(U_0)$ pour $\sigma\in N_H\setminus T_H\inserted{=N_H^-}$ \inserted{sont} tous égaux à un \unsure{m\textsuperscript{me}|même} ss-groupe $U_1$ de $H$

c) \deleted{$T_H$ commutatif, et $N_H/T_H$ \unsure{opère} par inversion sur $T_H$ par $t\mapsto t^{-1}$} $\sigma\in N_H^-\Longrightarrow\sigma^2=1$ ($\Rightarrow T_H$ est comm et $N_H/H\simeq\pm1$ y opère par $t\mapsto t^{-1}$)

d) $U_0\cap U_1=\{1\}$ (où $U_1$ défini dans b)).

On aimerait vérifier aussi la formule familière

\[
B_H\cap N_H=T_H\tag*{(164)}
\]

% slanted
\margin{donc la donnée de $(U_0,T_H,N_H)$ équivaut à celle de $(U_0,B_H,N_H)$, d'où on tire en effet $T_H=B_H\cap N_H$ ; \deleted{\unsure{puis}} \deleted{\unsure{et}} Dans certains cas importants, on aura d'ailleurs $U_0=[B_H,B_H]$ (en tous cas $\supset$) et tout est déterminé en termes de $B_H,N_H$ ss-groupes de $H$.}

\unsure{alors que|i.e.} le groupe $B_H\cap N_H$ est compris entre $T_H$ et $N_H$ -- donc si $\sigma\in N_H^-$ est donné, la formule (164) équivaut à

\[
\sigma\notin B_H\tag*{(165)}
\]

Mais si on avait $\sigma\in B_H\deleted{\unsure{=T_H.U_0=T_H|=B_H\,U_0\,T_H}}$, on aurait $\sigma\deleted{\unsure{\in}}(U_0)=\unsure{U_0|U_1}$ \inserted{de par d) $U_0=\{1\}$, donc $B_H=T_H$ donc $\sigma\in T_H$,} absurde.

\deleted{$u\in U_0,\ t\in T_H$, donc (comme $t$ centralise $T_H$, alors que $\sigma=ut$ le renverse et y induit $g\mapsto g^{-1}$) il faudrait que : $u\in U_0$ normalise $T_H$, et y induit $g\mapsto g^{-1}$ comme $T_H$ normalise $U_0$ et $T_H\cap U_0=\{1\}$, cel n'est possible que si $g\in T_H\Rightarrow g^2=1$ i.e \deleted{\unsure{de sorte}} $g=g^{-1}$ ($u$ \uline{centralise} $T_H$). \deleted{et} Donc $B_H\cap N_H=N_H$ i.e. $N_H\subset B_H$, \inserted{et alors}}

\[
\deleted{\left.\begin{aligned}
&\text{On a ou bien (164), \uline{ou} \uline{bien} }\quad T_H^2=\{1\}\ \text{et}\\
&N_H=\{1,u_0\}\times T_H,\ \text{où }u_0\in U_0\ \text{est un élément d'ordre }2\\
&\hspace{5em}\text{qui centralise }T_H,\ \text{i.e. }u_0\in{}_2(U_0^{T_H}),u_0\ne1.
\end{aligned}\right.}\tag*{\deleted{(165)}}
\]
% archive sheet number: 269
